%======================================================================
%  Memory subsystem: HBM2e, DDR5, memory modes and clustering
%======================================================================

\section{The memory subsystem}
\label{sec:xeon-max-platform}

Each Max 9480 socket has two local memory systems.  Eight DDR5 channels
attach conventional DIMMs, and four HBM2e stacks sit inside the package.
On the CRESCO8 HBM nodes every DDR5 channel holds one 64\,GB DIMM: the
EDAC labels list one DIMM on each of the eight channels of each socket,
so all channels are populated at one DIMM per channel and may run at the
full DDR5-4800 rate.  A socket therefore has 512\,GB of DDR5 behind
64\,GB of HBM, a capacity ratio of 8:1.  The idea of a small,
fast, on-package memory in front of a large DDR tier is not new; the
Knights Landing processor offered the same Flat and Cache choice for its
MCDRAM~\cite{peng2017hybridmemory}.  What is new for inference is that
the tier is attached to cores with AMX.

\subsection{Two local tiers and three access paths}
\label{sec:memory-hierarchy}

HBM and DDR5 trade different resources.  HBM offers several times the
bandwidth of DDR5 (\tabref{tab:xeon-max-peaks}); DDR5 offers eight times
the capacity and a lower unloaded latency.  HBM is therefore not simply
``faster memory''.  Its bandwidth is useful only when two conditions
hold together: the workload issues enough concurrent requests to
approach that bandwidth, and the bandwidth-critical data is local to the
cores that read it.

Every result must therefore name the access path it exercised:

\begin{itemize}
  \item \textbf{local HBM or local DDR5}, reached from cores on the same
        socket;
  \item \textbf{remote memory}, reached through UPI from the other
        socket; and
  \item \textbf{mixed placement}, in which pages, threads or model
        objects span several locality domains.
\end{itemize}

An apparently node-wide 128\,GB of HBM is consequently two 64\,GB local
tiers.  A workload uses the aggregate capacity and bandwidth only if its
data and threads are balanced across the sockets.
\tabref{tab:storage-levels} summarises the levels visible to one socket
and, for each, what decides its contents.

\begin{table}[htbp]
\centering
\small
\tabsetup
\begin{tabularx}{\textwidth}{P{2.3cm}P{3.3cm}P{3.3cm}Y}
\headrow
\headcell{Level} & \headcell{Capacity} & \headcell{Contents decided by} &
\headcell{Role in inference} \\[3pt]
\tabheadrule
\tabstripeRow
L1 data & 48\,KB per core & Hardware & Operands of the inner kernel
  loop. \\[3pt]

\tabrowsep
\tabplainRow
L2 & 2\,MB per core & Hardware & Kernel tiles and activations. \\[3pt]

\tabrowsep
\tabstripeRow
Last-level cache & 112.5\,MB, shared & Hardware & Smaller than any
  model's weights; reuse within a step. \\[3pt]

\tabrowsep
\tabplainRow
HBM2e & 64\,GB, four stacks & Software in Flat and HBM-only mode;
  hardware in Cache mode & Bandwidth tier. \\[3pt]

\tabrowsep
\tabstripeRow
DDR5 & 512\,GB, eight 64\,GB DIMMs on CRESCO8 & Software & Capacity
  tier. \\[3pt]

\tabrowsep
\tabplainRow
Remote memory & The other socket's tiers & Software, through thread
  and page binding & Reached over UPI; avoided by binding. \\[3pt]
\tabbottomrule
\end{tabularx}
\caption{Storage levels visible to one Max 9480 socket on a CRESCO8 HBM
node.  Cache capacities are those reported by \texttt{lscpu} on the
allocated node~\cite{intel9480}.}
\label{tab:storage-levels}
\end{table}

Published STREAM measurements show both the size of the bandwidth gap
and how much it varies between systems.  Na et al.\ report 588\,GB/s for
HBM and 233.8\,GB/s for DDR5 on a Xeon Max 9468
socket~\cite{na2024llmcpu}; Ibeid et al.\ report 680--840\,GB/s and
235--245\,GB/s on a Xeon Max 9470C~\cite{ibeid2025aurora}.  Neither
processor is the 9480, and the HBM figures differ by more than 15\%
between them.  The bandwidth of the CRESCO8 nodes is therefore measured
directly, in both Flat and Cache mode, in the next chapter.

\subsection{HBM memory modes}
\label{sec:hbm-modes}

The four HBM2e stacks are present in every package, but software does
not necessarily see them as memory.  The firmware's \emph{memory
hierarchy} setting gives HBM one of three roles: the only main memory, a
separately addressable tier beside DDR5, or a hardware-managed cache in
front of DDR5 (\figref{fig:hbm-modes}).  The setting is fixed at boot.
It is a property of the node, not a policy that a batch job can
select~\cite{lp1738,intelmaxtuning}.

\begin{figure}[htbp]
\centering
\begingroup
\definecolor{modeink}{HTML}{263445}
\definecolor{modeblue}{HTML}{365D85}
\colorlet{modeaccent}{sapred}
\definecolor{modeborder}{HTML}{CBD5DF}
\resizebox{\textwidth}{!}{%
\begin{tikzpicture}[
  font=\small\sffamily, text=modeink,
  card/.style={line width=0.65pt, rounded corners=4pt, align=center,
               inner xsep=5pt, inner ysep=7pt},
  panel/.style={card, draw=modeborder, fill=white,
                minimum width=4.35cm, minimum height=5.15cm},
  cpu/.style={card, draw=modeborder, fill=modeink!3,
              minimum width=3.7cm, minimum height=0.85cm},
  hbm/.style={card, draw=modeaccent, fill=modeaccent!9,
              text width=3.3cm, minimum width=3.7cm,
              minimum height=1.2cm},
  ddr/.style={hbm, draw=modeblue, fill=modeblue!9},
  absent/.style={card, draw=modeborder, densely dashed,
                 text=modeink!55, fill=modeink!1,
                 text width=3.3cm, minimum width=3.7cm,
                 minimum height=1.2cm},
  summary/.style={card, draw=modeborder, fill=modeink!3,
                  font=\footnotesize\sffamily, text width=3.95cm,
                  minimum width=4.35cm, minimum height=1.65cm},
  link/.style={{Stealth[length=5pt,width=4pt]}-{Stealth[length=5pt,width=4pt]},
               line width=0.8pt, draw=modeink!75},
  note/.style={font=\footnotesize\sffamily, text=modeink!80, align=center}
]
  \foreach \x in {0,4.85,9.7} {
    \node[panel] at (\x,0.6) {};
  }
  \node[font=\small\sffamily\bfseries] at (0,3.95) {HBM-only};
  \node[font=\small\sffamily\bfseries] at (4.85,3.95) {Flat (1LM)};
  \node[font=\small\sffamily\bfseries] at (9.7,3.95) {Cache (2LM)};
  \node[note] at (0,3.5) {HBM is the only main memory};
  \node[note] at (4.85,3.5) {Two separate main-memory tiers};
  \node[note] at (9.7,3.5) {HBM is a cache for DDR5};

  \node[cpu] (c0) at (0,2.5) {\textbf{CPU cores}};
  \node[cpu] (c1) at (4.85,2.5) {\textbf{CPU cores}};
  \node[cpu] (c2) at (9.7,2.5) {\textbf{CPU cores}};

  % HBM-only: one connected, allocatable tier.  DDR5 is not used.
  \node[hbm] (h0) at (0,0.7)
    {\textbf{HBM}\\{\footnotesize 64\,GB main memory}};
  \node[absent] at (0,-1.0)
    {\textbf{DDR5}\\{\footnotesize not used / not visible}};
  \draw[link, draw=modeaccent] (c0.south) -- (h0.north);

  % Flat: independent, parallel paths to two allocatable NUMA regions.
  % Side-by-side placement avoids suggesting that DDR5 is behind HBM.
  \node[hbm, text width=1.55cm, minimum width=1.85cm] (h1) at (3.8,0.7)
    {\textbf{HBM}\\{\footnotesize 64\,GB}};
  \node[ddr, text width=1.55cm, minimum width=1.85cm] (d1) at (5.9,0.7)
    {\textbf{DDR5}\\{\footnotesize RAM}};
  \draw[link, draw=modeaccent] ($(c1.south)+(-1.05,0)$) -- (h1.north);
  \draw[link, draw=modeblue] ($(c1.south)+(1.05,0)$) -- (d1.north);
  \node[note, text width=3.7cm] at (4.85,-1.0)
    {Separate NUMA regions.\\Software chooses the tier\\for each allocation.};

  % Cache: HBM is a transparent memory-side cache, not a RAM allocation target.
  \node[hbm, densely dashed] (h2) at (9.7,0.7)
    {\textbf{HBM cache}\\{\footnotesize 64\,GB; not allocatable}};
  \node[ddr] (d2) at (9.7,-1.0)
    {\textbf{DDR5}\\{\footnotesize main memory}};
  \draw[link, draw=modeaccent] (c2.south) -- (h2.north)
    node[midway, right, note, xshift=3pt] {lookup};
  \draw[link, draw=modeblue] (h2.south) -- (d2.north)
    node[midway, right, note, xshift=3pt] {miss / fill};

  % Identical summary rows make visibility, capacity and control comparable.
  \node[summary] at (0,-3.0)
    {\textbf{OS sees:} HBM only\\
     \textbf{RAM:} 64\,GB\\
     \textbf{Placement:} HBM only};
  \node[summary] at (4.85,-3.0)
    {\textbf{OS sees:} HBM + DDR5\\
     \textbf{RAM:} 64\,GB + DDR5\\
     \textbf{Placement:} software};
  \node[summary] at (9.7,-3.0)
    {\textbf{OS sees:} DDR5 only\\
     \textbf{RAM:} DDR5 capacity\\
     \textbf{Placement:} hardware};
\end{tikzpicture}%
}
\endgroup
\caption{The three Xeon Max memory modes, shown for one socket.  Solid
memory cards denote allocatable RAM; the dashed HBM card in Cache mode is
a transparent, hardware-managed cache whose capacity is not added to RAM.
Capacities are nominal; the paths are conceptual and omit the CPU's
internal caches~\cite{lp1603,lp1738,intelmaxtuning}.}
\label{fig:hbm-modes}
\end{figure}

\begin{description}
  \item[HBM-only mode.] HBM is the only main memory and DDR5 is absent
  from the operating system's address space.  Lenovo documents this
  topology for nodes without installed DDR5~\cite{lp1738}.  Applications
  need no placement mechanism, but the operating system, services and
  jobs share 64\,GB per socket.  It is the simplest route to HBM
  bandwidth when the whole workload fits.

  \item[Flat mode (1LM).] HBM and DDR5 are separate NUMA nodes.  Both
  capacities are available, and software can put a bandwidth-critical
  object in HBM while leaving colder data in DDR5.  That control is also
  an obligation: threads and pages must be bound with \texttt{numactl},
  libnuma or \texttt{mbind}, and actual residency must be verified,
  because first touch, page migration and file-backed mappings can place
  pages elsewhere.  Exposing the HBM nodes on Linux needs a second,
  per-boot administrative step that converts the HBM DAX devices to
  system RAM (\appref{app:flat-mode-setup}).

  \item[Cache mode (2LM).] Software sees only DDR5; HBM acts as a
  memory-side cache in front of it.  Unmodified programs benefit and the
  full DDR5 capacity is usable, but software cannot reserve HBM for
  particular data.  This is the mode of the regular CRESCO8 HBM
  partitions, and \secref{sec:cache-mode} describes it in detail.
\end{description}

The choice follows the working set and the control the software has
over placement.  HBM-only removes tier selection when everything fits;
Flat gives an HBM-aware program deterministic control over the hot set;
Cache gives the same large capacity with a transparent policy.  For
comparing the two tiers, Flat mode is the strongest design, because local
HBM and local DDR5 can be compared on the same socket without changing
the processor or the executable.  A Flat-versus-Cache comparison answers
a different question: explicit placement against a hardware cache policy.

\subsection{Cache mode in detail}
\label{sec:cache-mode}

The firmware describes the cache to the operating system through the
ACPI Heterogeneous Memory Attribute Table, and Linux exposes the
description in sysfs~\cite{linuxnumaperf}.  On the allocated CRESCO8 node
each DDR5 NUMA node reports a memory-side cache of 64\,GiB with 64-byte
lines, direct-mapped and write-back, and \texttt{lstopo} shows a 64\,GB
\texttt{MemCache} object above each 503\,GB NUMA node.  Each socket's
cache holds lines of its own socket's DDR5 only.  A thread that reads the
other socket's memory crosses UPI and is served by that socket's cache.

The mechanics follow from those attributes.  A read that hits is served
from HBM.  A miss reads the line from DDR5, returns it to the core, and
installs it in HBM; if the evicted line is dirty, it is first written
back to DDR5~\cite{intelmaxoverview}.  With 512\,GB of DDR5 behind 64\,GB
of cache, eight DDR5 lines compete for every cache slot, and because the
cache is direct-mapped each line has exactly one slot, selected by its
physical address.  Two consequences matter for measurement:

\begin{itemize}
  \item \textbf{A working set larger than the cache can be worse than no
        cache.}  Streaming beyond the capacity turns most accesses into
        misses, and every miss adds an HBM fill and possibly a write-back
        to the DDR5 traffic.  Effective bandwidth is therefore a property
        of the working-set size, and the next chapter measures it as
        such.
  \item \textbf{Fitting is a matter of physical pages, not only of
        size.}  Two hot pages whose physical addresses map to the same
        slots evict each other even when the total working set is below
        64\,GB.  The physical pages the kernel hands out vary from run to
        run, so conflict misses add run-to-run variance.  Measurements
        therefore need a warm-up pass, repetition across separate
        allocations, and a hot set comfortably below the capacity.
\end{itemize}

Cache mode does not prevent studying HBM for inference; it changes what
can be compared.  An 8-billion-parameter model in BF16 has about 16\,GB of
weights.  After warm-up, decode can stream those weights from the HBM
cache, so the same program run with a footprint below and above 64\,GB
per socket provides an HBM-versus-DDR contrast on any node of the
regular partitions.  What Cache mode cannot provide is object-level
control: weights or KV blocks cannot be pinned in HBM, cold data cannot
be excluded, and there is no DDR-only run with the same footprint on the
same node.  Those contrasts require the Flat-mode node.

The cache's hit rate is not directly observable by an ordinary user.  The
nodes run with \texttt{perf\_event\_paranoid} set to 2, which restricts
unprivileged users to their own processes' core events, and the uncore
memory-controller counters that would distinguish HBM from DDR5 traffic
are therefore unavailable (\secref{sec:cresco8-observed}).  Unless the
administrators grant counter access, residency in the cache has to be
inferred from working-set sweeps such as those of the next chapter.

\subsection{Clustering and NUMA}
\label{sec:hbm-clustering}

Memory mode is only the first configuration axis.  An independent
firmware setting decides how finely each socket is partitioned.  With
sub-NUMA clustering disabled, a setting Lenovo calls \emph{Quadrant},
the socket is one locality domain and addresses are interleaved across
all four tiles.  SNC4 exposes each tile, with its 14 cores, its share of
the last-level cache and its two DDR5 channels, as a separate NUMA domain
(\figref{fig:node-locality}).  When HBM is visible as memory, each tile's
16\,GB stack also becomes a separate node: in Flat mode these are
additional memory-only nodes, whereas in Cache mode they remain hidden.
SNC4 can raise local bandwidth for a NUMA-aware program, but it reduces
the largest local HBM allocation from 64\,GB to 16\,GB and creates more
ways to access memory remotely~\cite{lp1738,intelmaxtuning}.

Crossing three memory modes with two clustering modes gives the six
configurations in \tabref{tab:hbm-six-configurations}.  The NUMA counts
are for a complete dual-socket node, the unit allocated on CRESCO8, and
describe Flat mode after its DAX devices have been converted to system
RAM.  The axes are independent: every memory mode can be paired with
either clustering mode, so an experiment must report both.

\begin{table}[htbp]
\centering
\small
\tabsetup
\begin{tabularx}{\textwidth}{P{1.8cm}P{1.95cm}P{4.0cm}Y}
\headrow
\headcell{Memory mode} & \headcell{Cluster\-ing} &
\headcell{Visible topology (two sockets)} &
\headcell{Operational consequence} \\[3pt]
\tabheadrule
\tabstripeRow
HBM-only & Quadrant
  & 2 HBM nodes: one 64\,GB node per socket
  & Bind work per socket; all software consumes HBM. \\[3pt]

\tabrowsep
\tabplainRow
HBM-only & SNC4
  & 8 HBM nodes: four 16\,GB nodes per socket
  & Bind each worker and its allocation to the same tile. \\[3pt]

\tabrowsep
\tabstripeRow
Flat & Quadrant
  & 4 nodes: two DDR5 nodes with CPUs and two 64\,GB HBM memory-only nodes
  & Select the tier explicitly and pair each HBM node with cores on its
    local socket.  \textbf{Available on one CRESCO8 node.} \\[3pt]

\tabrowsep
\tabplainRow
Flat & SNC4
  & 16 nodes: eight DDR5 nodes with CPUs and eight 16\,GB HBM memory-only
    nodes
  & Select the tier and the matching tile; a misplaced page can be local
    to the socket but remote to the worker. \\[3pt]

\tabrowsep
\tabstripeRow
Cache & Quadrant
  & 2 DDR5 nodes, one per socket; HBM is hidden
  & Bind by socket; HBM contents are hardware-managed.  \textbf{Regular
    CRESCO8 HBM partitions.} \\[3pt]

\tabrowsep
\tabplainRow
Cache & SNC4
  & 8 DDR5 nodes, four per socket; HBM is hidden
  & Bind by tile; software still cannot choose what HBM holds. \\[3pt]
\tabbottomrule
\end{tabularx}
\caption{The six Xeon Max memory configurations and their OS-visible
NUMA topology on a dual-socket node.  ``Memory-only'' means a NUMA node
with capacity but no CPUs.  Node identifiers are assigned at boot and
must be read from the running system.  Synthesised from Lenovo's mode
and topology diagrams~\cite{lp1738}.}
\label{tab:hbm-six-configurations}
\end{table}

Flat mode is unforgiving because the HBM nodes it creates have no CPUs.
A default allocation by a thread on a CPU-bearing node lands in DDR5, and
selecting an HBM node without selecting the cores nearest to it creates
remote traffic.  A valid Flat experiment reads \texttt{numactl
-{}-hardware}, derives the CPU--HBM pairing from the distance matrix, binds
both threads and memory, and then verifies page residency.  The batch
system must also permit the binding.  In the regular CRESCO8 partitions
the cgroup of a job contains every NUMA node of the machine:
\texttt{cpuset.mems.effective} was \texttt{0-1} both for an exclusive
node and for an earlier 56-core allocation.  The Flat-mode node is not
allocated through these partitions (\secref{sec:cresco8-observed}), so
the question of whether a Slurm cpuset would admit its memory-only HBM
nodes does not arise for the experiments.

Lenovo's STREAM tests on an SD650~V3 with Max 9480 processors show how
large the binding effect is: in Flat and Quadrant mode, explicit binding
doubled the bandwidth of the unbound run~\cite{lp1738}.  The best
clustering mode also depends on the software.  Ibeid et al.\ obtain their
highest STREAM bandwidth with SNC4 and explicit locality, whereas Na et
al.\ observe worse inference performance and more remote traffic under
SNC4~\cite{ibeid2025aurora,na2024llmcpu}.
